% Source: https://artofproblemsolving.com/community/c3h2677201_a_problem_concept_ive_had_for_a_long_time
% Stable identity: aops-topic-2677201-q01
A rookable set is the graph of a permutation $\sigma$ of $\{0,1,2,3\}$ (square $(i,\sigma(i))$ in row $i$), so there are $24$ sets. Count orbits under the $8$ symmetries of the square with Burnside's lemma, by counting the sets each symmetry fixes:
\begin{itemize}
\item identity: $24$;
\item rotation by $180^\circ$, $(i,j)\mapsto(3-i,3-j)$: we need $\sigma(3-i)=3-\sigma(i)$, i.e. $\sigma$ maps the pairs $\{0,3\},\{1,2\}$ to pairs; $2\cdot2\cdot2=8$;
\item rotations by $\pm90^\circ$, $(i,j)\mapsto(j,3-i)$: we need $\sigma(\sigma(i))=3-i$, which forces the $4$-cycles $(0\,1\,3\,2)$ and $(0\,2\,3\,1)$; $2$ each;
\item reflections in the two mid-lines, e.g. $(i,j)\mapsto(3-i,j)$: we would need $\sigma(3-i)=\sigma(i)$, impossible; $0$ each;
\item reflections in the diagonals: $\sigma=\sigma^{-1}$ (or its analogue for the other diagonal), the $10$ involutions; $10$ each.
\end{itemize}
The number of orbits is $\dfrac{24+8+2+2+0+0+10+10}{8}=\boxed{7}$.
